%%%%%%%%%%%%%%%%%%%%%%%%%%%%%%%%%%%%%%%%%%%%%%%%
% Acrylverkauf: Plattengrößen und Preisfaktoren (Aushang)
% CC-BY-SA 3.0, FAU FabLab
% Früher Acrylpreis_Alt.pdf, jetzt LaTeX.
%%%%%%%%%%%%%%%%%%%%%%%%%%%%%%%%%%%%%%%%%%%%%%%%
\documentclass[a4paper,11pt]{scrartcl}
\usepackage[T1]{fontenc}
\usepackage[utf8]{inputenc}
\usepackage{lmodern}
\renewcommand{\familydefault}{\sfdefault}
\usepackage[ngerman]{babel}
\usepackage[a4paper,left=15mm,right=15mm,top=12mm,bottom=18mm,footskip=8mm]{geometry}
\usepackage{graphicx}
\usepackage[table]{xcolor}
\usepackage{tabularx}
\usepackage{tikz}
\usepackage{fancyhdr}

% Versionsnummer wie in fablab-document (version.tex erzeugt make)
\InputIfFileExists{version}{}{\newcommand{\Version}{Entwurf \today}}

\definecolor{fablab}{RGB}{20,51,105}
\definecolor{kopf}{gray}{0.88}

\pagestyle{fancy}
\fancyhf{}
\renewcommand{\headrulewidth}{0pt}
\fancyfoot[L]{\scriptsize Quelle: \texttt{github.com/fau-fablab/docs}}
\fancyfoot[R]{\scriptsize Version \Version}
\setlength{\parindent}{0pt}

% Eine Platte im Maßstab 1:6:  \Platte{x}{y}{Breite}{Höhe}{Name}{Preisfaktor}{Farbanteil}
% (Koordinaten und Maße in mm der echten Platte)
\newcommand{\Platte}[7]{%
  \fill[fablab!#7] (#1/6,#2/6) rectangle ++(#3/6,#4/6);
  \draw[black, line width=1.2pt] (#1/6,#2/6) rectangle ++(#3/6,#4/6);
  \node[align=center] at (#1/6+#3/12,#2/6+#4/12)
    {\textbf{#5}\\[0.5mm] \small #3\,×\,#4\,mm\\[0.5mm] \small Faktor \textbf{#6}};}

\begin{document}

% ------------------------------ Kopf ---------------------------------
\begin{minipage}[b]{0.62\linewidth}
  {\LARGE\bfseries Acrylverkauf}\\[2mm]
  \normalsize Acrylglas wird \textbf{nur} in \textbf{viertel, halben und ganzen Platten} verkauft.
\end{minipage}\hfill
\begin{minipage}[b]{0.35\linewidth}\raggedleft
  \IfFileExists{fablab-document/logo/Logo/Logo-FAU-bunt.pdf}%
    {\includegraphics[height=11mm]{fablab-document/logo/Logo/Logo-FAU-bunt.pdf}}{\textbf{FAU FabLab}}
\end{minipage}\\[1mm]
{\color{fablab}\rule{\linewidth}{0.8pt}}

\vspace{8mm}
{\Large\bfseries Wie funktioniert das genau?}\\[2mm]
Es gibt folgende Plattengrößen zu kaufen:\\[3mm]

\setlength{\tabcolsep}{6pt}
\renewcommand{\arraystretch}{1.5}
\arrayrulecolor{black!60}
\begin{tabularx}{\linewidth}{!{\color{black}\vrule width 1.2pt}>{\bfseries}p{0.34\linewidth}|p{0.24\linewidth}|X!{\color{black}\vrule width 1.2pt}}
  \noalign{\hrule height 1.2pt}
  \rowcolor{kopf} Name & \textbf{Maße} & \textbf{Preisfaktor} \\
  \noalign{\hrule height 0.8pt}
  ganze Platte        & 600\,×\,300\,mm & \textbf{1} (Grundpreis) \\ \hline
  \rowcolor{black!7}
  quadratische Platte & 300\,×\,300\,mm & \textbf{½} (halber Grundpreis) \\ \hline
  lange Platte        & 600\,×\,150\,mm & \textbf{½} (halber Grundpreis) \\ \hline
  \rowcolor{black!7}
  viertel Platte      & 300\,×\,150\,mm & \textbf{¼} (viertel Grundpreis) \\
  \noalign{\hrule height 1.2pt}
\end{tabularx}

\vspace{6mm}
Wenn du z.\,B. eine viertel Platte haben willst, schau zuerst, ob es schon eine vorgeschnittene gibt.
Wenn nicht, dann schneide erst eine ganze (oder halbe) Platte in vier (oder zwei) Teile.
Hierfür müssen natürlich \textbf{keine Lasergebühren} gezahlt werden.

\vspace{2mm}
Eine Schneidevorlage findest du in der Dateifreigabe unter \texttt{0\_Laservorlage}.

% ----------------------- Plattengrößen 1:6 ---------------------------
\vfill
\begin{center}
\begin{tikzpicture}[x=1mm,y=1mm]
  \Platte{0}{190}{600}{300}{ganze Platte}{1}{25}
  \Platte{650}{190}{300}{300}{quadratische Platte}{½}{15}
  \Platte{0}{0}{600}{150}{lange Platte}{½}{15}
  \Platte{650}{0}{300}{150}{viertel Platte}{¼}{8}
\end{tikzpicture}\\[3mm]
\small Plattengrößen im Maßstab 1\,:\,6
\end{center}
\vfill

\end{document}
